\begin{enumerate}

\item Obtener área, centroide, I\_z, I\_y e I\_yz sumando aportaciones de pared y booms.

\item Determinar los gradientes de tensión axial con ambos cortantes y el producto de inercia; integrar dq/ds en cada pared.

\item Imponer continuidad y equilibrio en los nudos; en contornos abiertos q=0 en extremos libres y en celdas cerradas imponer compatibilidad de giros.

\item Localizar el centro de cortadura anulando el torsor debido al flujo; calcular la tensión máxima por |q(s)|/t en paredes y puntos críticos.

\item Las leyes completas q(s) por pared se indican a continuación; la coordenada s se mide desde el nudo inicial de cada pared y el sentido firmado sigue las flechas del diagrama. Los resultados geométricos se obtienen con la misma idealización de paredes y larguerillos del enunciado.

\item Pared ala\_1: q(s) = $\frac{39574 s^{2} + 3088590 s}{12082785}$ N/mm; s en mm desde su nudo inicial, con signo según el recorrido.

\item Pared alma: q(s) = $\frac{- 540909 s^{2} + 57474780 s}{88607090} + \frac{21852120}{805519}$ N/mm; s en mm desde su nudo inicial, con signo según el recorrido.

\item Pared ala\_2: q(s) = $\frac{39574 s^{2} - 8389800 s}{12082785} + \frac{19826910}{805519}$ N/mm; s en mm desde su nudo inicial, con signo según el recorrido.

\item Pared crítica: alma; s = 53.1279568 mm desde su nudo inicial.

\end{enumerate}

\noindent\textbf{Operaciones y resultados intermedios.}

\begin{enumerate}

\item Área equivalente que soporta tensión axial.

\[\fitmath{y_{1}=2\cdot 60+2\cdot 110+2\cdot 45\approx 430\,\mathrm{mm^{2}}}\]

\item Propiedades geométricas de la sección (centroide y momentos de inercia).

\[\fitmath{z_G=-3.6627907\,\mathrm{mm},\quad y_G=51.1627907\,\mathrm{mm}}\]

\item Momentos de inercia del modelo en los ejes centroidales.

\[\fitmath{I_z=850751.938\,\mathrm{mm}^{4},\quad I_y=198981.105\,\mathrm{mm}^{4},\quad I_{yz}=303331.395\,\mathrm{mm}^{4}}\]

\item Centro de cortadura a partir del equilibrio de los flujos; coordenadas del modelo.

\[\fitmath{z_C=-4.31386473\,\mathrm{mm},\quad y_C=31.1077579\,\mathrm{mm}}\]

\item Tensión máxima |q|/t en pared alma.

\[\fitmath{y_{5}=\frac{\left|\frac{1.651933571\times10^{11}}{3724038263}\right|}{2}\approx 22.17933\,\mathrm{MPa}}\]

\end{enumerate}
